% DRAFT management-first abstract (not yet in main.tex). Numbers as in the
% audited draft; 3-seed and retraining results to be added when available.
\begin{abstract}
Autonomous management of a 5G cell carried by an unmanned aerial vehicle (UAV)
must keep search-and-rescue robots commandable without an operator and without
knowing the terrain. We built and flew, over an OpenAirInterface and Open5GS
testbed serving four robots, a three-loop management stack: a per-robot uplink
cap (every 0.2\,s), a liveness monitor (every 20\,s), and a terrain-blind
learning layer that forecasts link failure and places the UAV in three
dimensions (every 30\,s). The decisive loop was the simplest. Long forest
outages were not coverage loss but uplink queuing inside the user equipment
(UE): constant-rate video exceeded weak links and held command-and-control (C2)
traffic minutes behind it. The cap, driven only by round-trip time and bytes
sent, cut median C2 outage length from 143\,s to 17\,s and raised forest C2
availability from 66.8\,\% to 91.8\,\%. A second silent failure, dead UE
processes, had contaminated 12.4\,\% of training data. With both handled, the
learning layer forecasts C2 failure 20\,s ahead with modest skill that falls as
model capacity grows, does not recover how terrain changes the effect of
altitude, and in closed loop does not beat fixed hover or the robot centroid.
For autonomous management of such cells, sound inner loops (bounded queues,
verified liveness) must precede learned placement.
\end{abstract}
